\documentclass[11pt]{article}

\usepackage{fullpage}
\usepackage{amsmath}
\usepackage{epsfig}
\usepackage{graphics}
\usepackage{caption}
\usepackage{circuitikz}

\oddsidemargin -0.15in
\evensidemargin -0.15in
\headheight 0in
\headsep 0in
%\topmargin -0.25in
\textheight 9.3in
\textwidth 6.8in

\begin{document}


\begin{center}
\section*{
Franklin \& Marshall College \hfill PHY 223\\
\bigskip
{\Large \it Interferometry 2}}
\end{center}
 
In this second week or work with the photon interferometry apparatus, you will
observe and measure two-slit interference of photons with a source so weak that
there is only one photon in the apparatus at a time. You will need to refer to
the Week I lab handout as well as your notes and data from Week I.

\vspace{5mm}

\noindent {\bf Using the PMT}

In this lab, you will use the photomultiplier tube (PMT), which is a very
sensitive piece of equipment that can convert single photons into detectable
bursts of electrons. The PMT should not be exposed even to moderate levels of
light when turned off, and must not be exposed to anything but the dimmest of
lights when turned on. In this context, ordinary room light is intolerably
bright even to a PMT turned off, and light as dim as moonlight is much too
bright for a PMT turned on. That is why the PMT box is equipped with a shutter,
and the apparatus as a whole is equipped with a buzzer alarm, to help protect
the PMT from misadventure. The goal of these safety measures is to assure that
the PMT is used

\begin{itemize}
\item only when its box is coupled to the two-slit apparatus (it is), and
\item only when the cover of the apparatus is closed, and
\item only when the light-bulb (and not the laser) source is being used inside
the box.
\end{itemize}

Before you do anything with the PMT, locate on the detector box the HIGH-VOLTAGE
toggle switch that activates its power supply, and ensure that it’s turned off;
also turn down to 0.00 the 10-turn dial that will later set the voltage that
enables the process of electron-multiplication inside it. Also, ensure that the
shutter of the box is still in its down position.

\vspace{5mm}

\noindent {\bf Procedure: Single Photon Mode}

Move the laser out of the way so that light from the bulb can propagate through
the system. Turn the bulb power adjustment down to 0 on its scale, and set the
3-position toggle switch to the BULB position; dial the bulb adjustment up from 0 until
you see the bulb light up. Keep the amplitude to below 6 to help prevent the
bulb from burning out early. The light you observe should appear green. If it
doesn't, you're missing a filter that only lets light between 541 nm and 551 nm. 

Close the enclosure carefully. Adjust the micrometers so that both slits are
open and the detector slit is at the center of the interference pattern. Connect
the computer to the PMT box and start the software to measure the counts on the
PMT. {\bf Shut the door and turn the lights off.} Turn the PMT on and set the
gain to 4. Turn the bulb on and adjust its amplitude until the PMT registers
10$^3$-10$^4$ events per seconds.

Measure the number of events for each position of the detector slit across the
same range as your previous measurement with the photodiode. At each location,
let the computer update a number of times and note a couple of different values. 

Move the central micrometer to block both slits and record the dark counts. You
will need this for your analysis.

\vspace{5mm}

\noindent {\bf Analysis}

\begin{itemize}
\item Graph your data with uncertainty.
\item Compare your results to your model.
\item Estimate the photon density. To do this, note that the PMT is not 100\%
efficient; in fact, for green light, it produces countable electronic events for
only about 4\% of the incident photons. Now take your central-maximum reading
for the event rate, and use this efficiency estimate to infer the rate of
arrival of all photons at the photocathode of the PMT. From this arrival rate,
compute the average time interval between the arrivals of successive photons.
Next, compute the ‘time of flight’ of a photon in the apparatus, the time it
takes to traverse the distance from the bulb source to the detector. Comparing
the time of flight of a given photon to the typical time interval between the
arrival of successive photons, you can compute the fraction of the time there is
even one photon in flight through the apparatus. What is this fraction for your
data? This fraction measures numerically the sense in which the apparatus allows
you to work ‘one photon at a time’; and yet even at this low rate of photons,
you can see phenomena that you attribute to constructive and destructive
interference.
\end{itemize}
\vspace{5mm}

\noindent {\bf (Optional) Adjusting the PMT}

You’ll need a reasonable fast ($>$20 MHz bandwidth), preferable digital,
oscilloscope for first examination, and a digital counter sensitive to TTL-level
($+$4 V positive-going) pulses for eventual counting, of photon events. Set the
scope to about 50 mV/division vertically, and 200 ns/division horizontally, and
set it to trigger on positive-going pulses of perhaps $>$20 mV height. Now find
the PHOTOMULTIPLIER OUTPUT of the detector box, and connect it via a BNC cable
to the vertical input of the scope; use a 50 $\Omega$ termination at the scope.
You are about to look at pulses, each of which starts with the light-induced
ejection of a single electron from the light-sensitive photocathode of the PMT.
Those photoelectrons will be amplified by about 10$^5$ inside the PMT, and
arrive as a pulse of about 10$^5$ electrons, or about -2x10$^{-14}$ Coulombs,
at the input of a charge- sensitive amplifier. That device will convert the
pulse of negative charge to a positive- going voltage pulse, which you will see
on the scope.

When you are ready with these electronics, keep the shutter closed, set the
HIGH- VOLTAGE 10 - turn dial to 0.00, and turn on the HIGH-VOLTAGE toggle
switch. You are now ready to start dialing up the high-voltage supply that
provides the potentials inside the PMT needed for the electron-multiplication
process; watch the scope display as you start to dial up the voltage, about a
turn at a time. [Each full turn raises the PMT ‘bias’ by about 100 V, and the
``gain'' of the PMT grows exponentially with this voltage setting. A pair of
terminals on the panel allows you to monitor the potential difference (PMT
bias/1000) using an ordinary digital multimeter.] As you raise the bias voltage,
you will see occasional transient electronic noise, but the scope display should
be quiet at each steady voltage. Somewhere around a setting of 4 or 5 turns of
the dial, you should get occasional positive-going pulses on the scope,
occurring at a modest rate of 1-10 per second.

\begin{itemize}
\item If you see some sinusoidal modulation of a few mV amplitude, and of about
200 kHz frequency, in the baseline of the PMT signal, this is normal.

\item If you see a continuing high rate ($>$10 kHz) of pulses from the PMT, this
is not normal, and you should turn down, or off, the bias level and start
fresh - you may have a malfunction, or a light leak.
\end{itemize}

If you see this low rate of pulses, you have discovered the ``dark rate'' of the
PMT, its output pulse rate even in the total absence of light. You also now have
the PMT ready to look at photons from your two-slit apparatus, so finally you
may open the shutter (by grasping that vertically-emerging rod and lifting it a
few cm). The scope should now show a much greater rate of pulses, perhaps of
order 10$^3$ per second, and that rate should vary systematically with the setting
of the bulb intensity.

If you can see those analog pulses, you are ready to look at the digital pulses
they trigger. Inside the PMT pulse amplifier is a pulse generator that emits, at
the OUTPUT TTL connector, a single pulse, of fixed height and duration, each
time the analog pulse you’re viewing exceeds an adjustable threshold. Arrange
another BNC cable, also terminated at its far end, to a second vertical channel
on your scope, this one set for 2 V/div vertically. Now start with the
discriminator control on your PMT box set to 0, and watch the scope display the
analog, and TTL-level, pulse outputs in dual-trace operation. You should be able
to find a discriminator setting, low on the dial, for which the scope shows one
TTL pulse for each of, and for only, those analog pulses that reach (say) a $+$50
mV level.

\begin{itemize}
\item If your analog pulses are mostly not this high, you can raise the PMT bias
by half a turn (50 Volts) to gain more electron multiplication.

\item If your TTL pulses come much more frequently than the analog pulses, set
the discriminator dial lower on its scale.
\end{itemize}

Now send the TTL pulses to a counter, arranged to display successive readings of
the number of pulses that occur in successive 1 sec intervals. These represent
photon-induced events (which are converted in the PMT to photoelectrons, then
multiplied to electron pulses, then amplified to voltage pulses meeting a
threshold criterion). Because of the <100\% efficiency of the PMT, not every
photon is detected; but you are seeing a representative subset of events due to
the arrivals of individual photons.

To confirm that this is true, record a series of ``dark counts'' obtained with
the light bulb dialed all the way town to 0 on its scale. Now choose a setting
that gives an adequate photon count rate (about 10$^3$/sec) and use the
slit-blocker, according to your previously obtained settings, to block the light
from both slits. This should reduce the count rate to a background rate,
probably somewhat higher than the dark rate. Next, open up both slits, and try
moving the detector slit to see if you can see interference fringes in the
photon count rate. You will need to pick an detector-slit location, then wait
for a second or more, then read the photon count in one or more 1 sec intervals,
before trying a new detector-slit location. If you can see maxima and minima,
you are ready to take data.

A first sort of data will enable you to be assured that you have set the PMT
bias voltage to a suitable range. For this data, the independent variable is the
PMT voltage setting, or one of its surrogates the 10-turn dial setting (1.00
turn $=$ 100 V of bias) or the monitor voltage (0.10 V $=$ 100 V of bias). The
dependent variables are both count rates as read by the TTL counter: one is the
count rate at the central maximum of the interference pattern, with both slits
open; the other is the dark rate, the count rate obtained under identical
conditions but with the PMT shutter closed. Try recording both rates over the
range 300- 650 V of PMT bias, and plot the two count rates on a semi-logarithmic
graph. (Or you can plot the log of the count rate against the bias voltage.) You
should see the ``light rate'' reach a plateau, with the interpretation that you
have reached a PMT bias that suffices to allow each genuine photoelectron to
trigger the whole chain of electronics all the way to the TTL counter; you
should also see the (much lower) ``dark rate'' also rising with PMT bias. See if
your graph motivates a setting at which you are counting substantially all of
the honest-to-goodness photon events, but minimizing the number of ``dark
events''.

\end{document}